\documentclass[../../main.tex]{subfiles}
\begin{document}
\begin{oframed}
  \subfile{problem}
\end{oframed}

{\bf [解答]}

\paragraph{(1)}

$p + q\sqrt{2} + r\sqrt{3} = 0 \iff q\sqrt{2} + r\sqrt{3} = -p$ だから両辺2乗して
\begin{align}
  (2q^2 + 3r^2 - p^2) + 2qr\sqrt{6} = 0
\end{align}
である．$qr \neq 0$とすると
\begin{align}
  \sqrt{6} = -\dfrac{2q^2 + 3r^2 - p^2}{2qr}
\end{align}
となるが，$p, q, r \in \mathbb{Q}$ より右辺は有理数，左辺は無理数となり矛盾する．
従って
\begin{align}
  \begin{cases}
    2q^2 + 3r^2 - p^2 = 0  \\
    qr = 0
  \end{cases} \label{eq:1}
\end{align}
となる．第二式から $q=0$ 又は $r=0$ である．

\subparagraph{$1^\circ \; q=0$ の時}

\cref{eq:1}の第一式から $3r^2 = p^2$ である．
よって$p = \pm \sqrt{3} r$ となって，$r \ne 0$ の時 $\sqrt{3} = \pm\dfrac{p}{r}$ で右辺が有理数，左辺が無理数となり矛盾する．したがって
\begin{align}
  p = r = 0
\end{align}
である．

\subparagraph{$2^\circ \; r=0$ の時}

\cref{eq:1}の第一式から $2q^2 = p^2$ であり $p = \pm \sqrt{2} q$ となって，$1^{\circ}$と同様に $p = q = 0$が従う．
\casesend

以上からいずれの場合も $p = q = r = 0$ である．よって題意は示された．

\paragraph{(2)}

題意を背理法で示す．
$f(1) = a+b+1$, $f(1+\sqrt{2}) = (a+b+3) + (a+2)\sqrt{2}$, $f(\sqrt{3}) = 3+b + a\sqrt{3}$ がいずれも有理数と仮定すると，有理数 $\alpha, \beta, \gamma$ を用いて，
\begin{align}
  & a+b = \alpha  \label{eq:2}\\
  & (a+2)\sqrt{2} = \beta \label{eq:3} \\
  & b+a\sqrt{3} = \gamma \label{eq:4}
\end{align}
とかける．

また $f(x)$ の係数は実数ゆえ $a, b \in \mathbb{R}$ である．
\cref{eq:2,eq:4}から$b$を消去すると
\begin{align}
  a(1 - \sqrt{3}) = \alpha - \gamma
  \quad \therefore
  a = (\gamma - \alpha) \cdot \dfrac{1+\sqrt{3}}{2}
\end{align}
だから\cref{eq:3}に代入して
\begin{align}
  & \dfrac{1}{2} \left[ (\gamma - \alpha)(1 + \sqrt{3}) + 4 \right] \sqrt{2} = \beta \\
  \therefore
  & \dfrac{1}{2}(\gamma - \alpha + 4)\sqrt{2} - \beta + \dfrac{1}{2}(\gamma - \alpha)\sqrt{6} = 0
\end{align}
を得る． $\alpha, \beta, \gamma \in \mathbb{Q}$ 及び (1)から
\begin{align}
  \begin{cases}
    \gamma - \alpha + 4 = 0 \\
    \beta = 0 \\
    \gamma - \alpha = 0
  \end{cases}
\end{align}
となるが，第一式，第三式より $4=0$ となり矛盾する．
したがって背理法により題意は示された．

\end{document}
